% ============================================================
\chapter{Data Pipelines and Feature Stores}
\label{ch:pipelines}
\index{data pipelines}
% ============================================================

\section{Why Data Pipelines?}
\label{sec:why-pipelines}

In production ML, data rarely arrives clean and ready. A \textbf{data
pipeline} automates the sequence of steps that transform raw data into
features suitable for training or inference.

\begin{definition}[Data pipeline]
A \textbf{data pipeline}\index{data pipeline} is an automated,
orchestrated sequence of data processing steps (extraction, transformation,
validation, loading) that converts raw data sources into curated datasets
ready for analysis or model training.
\end{definition}

\begin{remark}
Without automated pipelines, data scientists spend up to 80\% of their
time on data wrangling---manually downloading, cleaning, and
transforming data for each experiment.
\end{remark}

\section{ETL vs ELT}
\label{sec:etl-elt}
\index{ETL}\index{ELT}

\begin{center}
\begin{tabular}{@{}lp{5.5cm}p{5.5cm}@{}}
\toprule
& \textbf{ETL} & \textbf{ELT} \\
\midrule
\textbf{Order} & Extract $\to$ Transform $\to$ Load &
  Extract $\to$ Load $\to$ Transform \\
\textbf{Where} & Transformation on an intermediate server &
  Transformation inside the data warehouse \\
\textbf{Best for} & Structured data, relational sources &
  Large-scale, cloud data warehouses \\
\textbf{Tools} & Airflow, Prefect, Luigi &
  dbt, BigQuery, Snowflake \\
\textbf{ML use case} & Feature engineering pipelines &
  Analytics features from data warehouse \\
\bottomrule
\end{tabular}
\end{center}

\section{Pipeline Architecture}
\label{sec:pipeline-architecture}

\begin{center}
\begin{tikzpicture}[
  box/.style={rectangle, draw, rounded corners, minimum width=2.5cm,
    minimum height=1.2cm, align=center, font=\small},
  storage/.style={cylinder, draw, shape border rotate=90,
    minimum width=1.8cm, minimum height=1cm, align=center,
    font=\scriptsize, aspect=0.3},
  arr/.style={-{Stealth[length=2.5mm]}, thick, steelblue}
]
  % Sources
  \node[storage, fill=blue!10] (src1) {\textbf{Database}\\PostgreSQL};
  \node[storage, fill=blue!10, below=0.8cm of src1] (src2) {\textbf{API}\\REST/gRPC};
  \node[storage, fill=blue!10, below=0.8cm of src2] (src3) {\textbf{Files}\\CSV, JSON};

  % Extract
  \node[box, fill=green!10, right=2cm of src2] (extract) {\textbf{Extract}\\Ingest};

  % Transform
  \node[box, fill=yellow!10, right=1.8cm of extract] (transform)
    {\textbf{Transform}\\Clean, Engineer};

  % Validate
  \node[box, fill=orange!10, right=1.8cm of transform] (validate)
    {\textbf{Validate}\\Schema, Quality};

  % Destinations
  \node[storage, fill=red!10, right=2cm of validate] (fs)
    {\textbf{Feature Store}\\Feast};
  \node[storage, fill=red!10, above=0.5cm of fs] (wh)
    {\textbf{Warehouse}\\BigQuery};
  \node[storage, fill=red!10, below=0.5cm of fs] (train)
    {\textbf{Training}\\Dataset};

  % Arrows
  \draw[arr] (src1.east) -- (extract.west);
  \draw[arr] (src2.east) -- (extract.west);
  \draw[arr] (src3.east) -- (extract.west);
  \draw[arr] (extract) -- (transform);
  \draw[arr] (transform) -- (validate);
  \draw[arr] (validate.east) -- (wh.west);
  \draw[arr] (validate.east) -- (fs.west);
  \draw[arr] (validate.east) -- (train.west);

  % Orchestrator
  \node[box, fill=gray!10, below=2cm of transform, minimum width=8cm]
    (orch) {\textbf{Orchestrator}: Airflow / Prefect / Dagster};
  \draw[dashed, gray, thick] (orch.north) -- (extract.south);
  \draw[dashed, gray, thick] (orch.north) -- (transform.south);
  \draw[dashed, gray, thick] (orch.north) -- (validate.south);
\end{tikzpicture}
\end{center}

\section{Orchestration with Airflow}
\label{sec:airflow}
\index{Airflow}

\textbf{Apache Airflow}\index{Airflow} is the most widely used workflow
orchestrator. It defines pipelines as Python DAGs (Directed Acyclic Graphs)
and provides a web UI for monitoring, scheduling, and retry management.

\begin{codeblock}[title=\textbf{Airflow DAG for an ML pipeline}]
\begin{minted}{python}
from airflow import DAG
from airflow.operators.python import PythonOperator
from datetime import datetime, timedelta

default_args = {
    "owner": "ml-team",
    "retries": 2,
    "retry_delay": timedelta(minutes=5),
}

with DAG(
    dag_id="ml_training_pipeline",
    default_args=default_args,
    schedule_interval="@daily",
    start_date=datetime(2025, 1, 1),
    catchup=False,
    tags=["ml", "training"],
) as dag:

    def extract_data(**kwargs):
        """Extract data from source."""
        import pandas as pd
        df = pd.read_sql("SELECT * FROM events", conn)
        df.to_parquet("/data/raw/events.parquet")

    def transform_data(**kwargs):
        """Clean and engineer features."""
        import pandas as pd
        df = pd.read_parquet("/data/raw/events.parquet")
        df = df.dropna(subset=["target"])
        df["hour"] = pd.to_datetime(df["ts"]).dt.hour
        df.to_parquet("/data/processed/features.parquet")

    def train_model(**kwargs):
        """Train and log the model."""
        import mlflow
        # ... training with MLflow logging ...

    extract = PythonOperator(
        task_id="extract_data", python_callable=extract_data
    )
    transform = PythonOperator(
        task_id="transform_data", python_callable=transform_data
    )
    train = PythonOperator(
        task_id="train_model", python_callable=train_model
    )

    extract >> transform >> train
\end{minted}
\end{codeblock}

\begin{warning}[title=\textbf{Airflow pitfalls}]
\begin{itemize}
  \item Airflow is an \textbf{orchestrator}, not a data processing engine.
    Heavy computations should run on Spark, Dask, or dedicated workers.
  \item The scheduler can be resource-hungry; plan infrastructure
    accordingly.
  \item DAG files are parsed frequently---keep them lightweight (no heavy
    imports at module level).
\end{itemize}
\end{warning}

\section{Orchestration with Prefect}
\label{sec:prefect}
\index{Prefect}

\textbf{Prefect}\index{Prefect} is a modern alternative to Airflow with a
simpler API. Flows are regular Python functions decorated with
\cli{@flow} and \cli{@task}.

\begin{codeblock}[title=\textbf{Complete Prefect ML pipeline}]
\begin{minted}{python}
from prefect import flow, task
from prefect.tasks import task_input_hash
from datetime import timedelta
import pandas as pd
from sklearn.ensemble import RandomForestClassifier
from sklearn.model_selection import train_test_split
from sklearn.metrics import accuracy_score
import mlflow

@task(retries=3, cache_key_fn=task_input_hash,
      cache_expiration=timedelta(hours=1))
def extract_data(source: str) -> pd.DataFrame:
    """Extract data from source with caching."""
    df = pd.read_csv(source)
    print(f"Extracted {len(df)} rows")
    return df

@task
def validate_data(df: pd.DataFrame) -> pd.DataFrame:
    """Validate data quality."""
    assert len(df) > 100, "Dataset too small"
    assert df.isnull().mean().max() < 0.3, "Too many nulls"
    null_pct = df.isnull().mean()
    print(f"Max null %: {null_pct.max():.2%}")
    return df.dropna()

@task
def engineer_features(df: pd.DataFrame) -> pd.DataFrame:
    """Create features for training."""
    df["feature_ratio"] = df["col_a"] / (df["col_b"] + 1)
    df["feature_log"] = df["col_c"].apply(
        lambda x: max(x, 0)
    ).apply(lambda x: x ** 0.5)
    return df

@task
def train_and_evaluate(
    df: pd.DataFrame,
    target_col: str = "target",
    n_estimators: int = 200,
) -> dict:
    """Train model and return metrics."""
    X = df.drop(columns=[target_col])
    y = df[target_col]
    X_train, X_test, y_train, y_test = train_test_split(
        X, y, test_size=0.2, random_state=42
    )

    with mlflow.start_run():
        model = RandomForestClassifier(
            n_estimators=n_estimators, random_state=42
        )
        model.fit(X_train, y_train)

        acc = accuracy_score(y_test, model.predict(X_test))
        mlflow.log_param("n_estimators", n_estimators)
        mlflow.log_metric("accuracy", acc)
        mlflow.sklearn.log_model(model, "model")

    return {"accuracy": acc, "n_samples": len(df)}

@flow(name="ML Training Pipeline",
      description="End-to-end training pipeline")
def training_pipeline(
    data_source: str = "data/dataset.csv",
    n_estimators: int = 200,
):
    """Main pipeline flow."""
    raw_data = extract_data(data_source)
    clean_data = validate_data(raw_data)
    features = engineer_features(clean_data)
    metrics = train_and_evaluate(
        features, n_estimators=n_estimators
    )
    print(f"Pipeline complete. Accuracy: {metrics['accuracy']:.4f}")
    return metrics

if __name__ == "__main__":
    training_pipeline()
\end{minted}
\end{codeblock}

\begin{codeblock}[title=\textbf{Running and scheduling Prefect flows}]
\begin{minted}{bash}
# Run locally
python pipeline.py

# Start the Prefect server
prefect server start

# Deploy with a schedule
prefect deployment build pipeline.py:training_pipeline \
    --name "daily-training" \
    --cron "0 6 * * *"
prefect deployment apply training_pipeline-deployment.yaml

# Start a worker to execute scheduled runs
prefect worker start -p default-agent-pool
\end{minted}
\end{codeblock}

\section{DVC Pipelines for ML}
\label{sec:dvc-pipelines-ch5}
\index{DVC!pipelines}

As seen in Chapter~\ref{ch:version-control}, DVC pipelines define
reproducible ML workflows in \file{dvc.yaml}. Unlike Airflow/Prefect,
DVC pipelines are \textbf{data-driven} (stages re-run only when inputs
change) rather than \textbf{schedule-driven}.

\begin{center}
\begin{tabular}{@{}lcc@{}}
\toprule
\textbf{Aspect} & \textbf{DVC Pipelines} & \textbf{Airflow/Prefect} \\
\midrule
Trigger         & Data/code changes & Schedule or event \\
State tracking  & File hashes       & Database \\
Caching         & Content-addressable & Time-based \\
Deployment      & CLI (\cli{dvc repro}) & Server + workers \\
Best for        & Experimentation   & Production orchestration \\
\bottomrule
\end{tabular}
\end{center}

\section{Feature Stores}
\label{sec:feature-stores}
\index{feature store}

A \textbf{feature store}\index{feature store} is a centralised repository
for storing, sharing, and serving ML features. It solves the problem of
feature reuse and training-serving skew.

\begin{definition}[Feature store]
A \textbf{feature store} is a data system that manages the lifecycle of
ML features: definition, computation, storage, versioning, and serving
for both training (batch) and inference (online, low-latency).
\end{definition}

\begin{warning}[title=\textbf{Training-serving skew}]
\textbf{Training-serving skew}\index{training-serving skew} occurs when
features are computed differently at training time and inference time.
For example, a feature computed as a 7-day rolling average in a batch
training script may be computed as a simple average in the serving code.
A feature store ensures the \emph{same} transformation is used in both
contexts.
\end{warning}

\subsection{Feast}
\index{Feast}

\textbf{Feast}\index{Feast} (Feature Store) is an open-source feature
store that provides offline (batch) and online (low-latency) feature
serving.

\begin{codeblock}[title=\textbf{Feast feature definitions}]
\begin{minted}{python}
# feature_repo/features.py
from feast import Entity, FeatureView, Field, FileSource
from feast.types import Float32, Int64
from datetime import timedelta

# Define the entity (the business object)
customer = Entity(
    name="customer_id",
    join_keys=["customer_id"],
    description="Unique customer identifier",
)

# Define the data source
customer_source = FileSource(
    path="data/customer_features.parquet",
    timestamp_field="event_timestamp",
)

# Define features
customer_features = FeatureView(
    name="customer_features",
    entities=[customer],
    ttl=timedelta(days=90),
    schema=[
        Field(name="total_purchases", dtype=Int64),
        Field(name="avg_order_value", dtype=Float32),
        Field(name="days_since_last_order", dtype=Int64),
        Field(name="purchase_frequency", dtype=Float32),
    ],
    source=customer_source,
    online=True,
    tags={"team": "ml-platform"},
)
\end{minted}
\end{codeblock}

\begin{codeblock}[title=\textbf{Using Feast for training and serving}]
\begin{minted}{python}
from feast import FeatureStore
import pandas as pd

store = FeatureStore(repo_path="feature_repo/")

# --- Offline: get historical features for training ---
entity_df = pd.DataFrame({
    "customer_id": [1001, 1002, 1003],
    "event_timestamp": pd.to_datetime(["2025-01-15"] * 3),
})

training_df = store.get_historical_features(
    entity_df=entity_df,
    features=[
        "customer_features:total_purchases",
        "customer_features:avg_order_value",
        "customer_features:days_since_last_order",
        "customer_features:purchase_frequency",
    ],
).to_df()

print(training_df.head())

# --- Online: get features for real-time inference ---
# First, materialise features to the online store
# feast materialize 2024-01-01T00:00:00 2025-06-01T00:00:00

online_features = store.get_online_features(
    features=[
        "customer_features:total_purchases",
        "customer_features:avg_order_value",
    ],
    entity_rows=[{"customer_id": 1001}],
).to_dict()

print(online_features)
\end{minted}
\end{codeblock}

\begin{codeblock}[title=\textbf{Feast CLI commands}]
\begin{minted}{bash}
# Initialise a feature repository
feast init feature_repo
cd feature_repo

# Apply feature definitions
feast apply

# Materialise features to the online store
feast materialize 2024-01-01T00:00:00 2025-06-01T00:00:00

# Serve features via REST API
feast serve
\end{minted}
\end{codeblock}

\section{Orchestrator Comparison}
\label{sec:orchestrator-comparison}

\begin{center}
\begin{tabular}{@{}lcccc@{}}
\toprule
\textbf{Criterion} & \textbf{Airflow} & \textbf{Prefect} & \textbf{Dagster} & \textbf{Luigi} \\
\midrule
Language       & Python  & Python  & Python  & Python \\
DAG definition & Python  & Python decorators & Python + assets & Python classes \\
UI             & Good    & Excellent & Excellent & Basic \\
Scheduling     & Cron, event & Cron, event & Cron, sensor & Cron \\
Caching        & No      & Yes    & Yes    & Yes    \\
Data lineage   & Limited & Good   & Excellent & Limited \\
Cloud managed  & Yes     & Yes    & Yes    & No     \\
Community      & Very large & Growing & Growing & Small \\
Learning curve & Medium  & Low    & Medium & Low    \\
\bottomrule
\end{tabular}
\end{center}

\begin{bestpractice}[title=\textbf{Choosing an orchestrator}]
\begin{itemize}
  \item \textbf{Airflow}: battle-tested, huge ecosystem, best for large
    teams with complex scheduling needs.
  \item \textbf{Prefect}: modern API, easy to start, best for teams that
    want Python-native pipelines without boilerplate.
  \item \textbf{Dagster}: asset-centric paradigm, excellent for data
    engineering teams that think in terms of data assets.
  \item \textbf{DVC}: best for ML experimentation where pipelines are
    re-run on data/code changes, not on a schedule.
\end{itemize}
\end{bestpractice}

\section{Best Practices}

\begin{bestpractice}[title=\textbf{Pipeline design golden rules}]
\begin{enumerate}
  \item \textbf{Idempotency}: running a pipeline twice with the same input
    must produce the same output.
  \item \textbf{Atomicity}: each step either fully succeeds or fully fails
    (no partial writes).
  \item \textbf{Data validation}: validate schemas and distributions at
    every stage boundary.
  \item \textbf{Retry logic}: transient failures (network, API rate limits)
    should be retried with backoff.
  \item \textbf{Logging and alerting}: every step should log its progress
    and alert on failure.
  \item \textbf{Separation of concerns}: orchestration logic (DAG
    definition) should be separate from business logic (transformations).
\end{enumerate}
\end{bestpractice}

\begin{antipattern}[title=\textbf{Pipeline anti-patterns}]
\begin{itemize}
  \item Running heavy computation inside the orchestrator process
  \item Hardcoding file paths and connection strings
  \item No data validation between pipeline stages
  \item Pipelines that cannot be re-run safely (non-idempotent writes)
  \item Feature computation duplicated between training and serving code
\end{itemize}
\end{antipattern}

\section{Mini-project: Orchestrated ML Pipeline}
\label{sec:mini-project-ch5}

\begin{miniproject}[title=\textbf{Mini-project 5: End-to-end orchestrated pipeline}]
Build a complete data-to-model pipeline:
\begin{enumerate}
  \item Write a Prefect flow with 4 tasks: \texttt{extract},
    \texttt{validate}, \texttt{engineer\_features}, \texttt{train}.
  \item Add data validation checks (schema, null percentage, class
    balance).
  \item Integrate MLflow tracking in the training task.
  \item Add retry logic and caching for the extraction task.
  \item Deploy the pipeline with a daily schedule.
  \item (Bonus) Define 2--3 features in Feast and serve them for
    training and online inference.
\end{enumerate}
\end{miniproject}

\section{Exercises}
\label{sec:exercises-ch5}

\begin{exercise}[$\bigstar$ --- Prefect basics]
Write a Prefect flow with three tasks: (1)~download a CSV from a URL,
(2)~compute basic statistics (mean, std, null count per column),
(3)~save a summary report to a JSON file. Add retry logic to the download
task. Run the flow and inspect the Prefect UI.
\end{exercise}

\begin{exercise}[$\bigstar\bigstar$ --- Pipeline with validation]
Build a data pipeline (Prefect or Airflow) for a tabular ML task that
includes:
\begin{enumerate}
  \item Data extraction from two sources (CSV + SQLite database)
  \item Data joining and cleaning
  \item Schema validation using \pkg{pandera} or \pkg{great-expectations}
  \item Feature engineering (at least 3 derived features)
  \item Train/test split with the split ratio logged
\end{enumerate}
Verify that the pipeline fails gracefully when given malformed data.
\end{exercise}

\begin{exercise}[$\bigstar\bigstar\bigstar$ --- Feature store integration]
Build a feature store using Feast for a customer churn prediction task:
\begin{enumerate}
  \item Define at least 5 features across 2 feature views
  \item Materialise features to the online store
  \item Write a training script that pulls historical features from Feast
  \item Write a serving script that pulls online features for real-time
    prediction
  \item Verify consistency between offline and online feature values
  \item Document the feature definitions and their business meaning
\end{enumerate}
\end{exercise}

\section{Cheatsheet}

\begin{cheatsheet}[title=\textbf{Chapter 5 Summary}]
\begin{tabular}{@{}p{4cm}p{8.5cm}@{}}
\toprule
\textbf{Tool / Concept} & \textbf{Key Points} \\
\midrule
\textbf{ETL vs ELT} & ETL transforms before loading; ELT transforms
  inside the warehouse \\
\textbf{Airflow} & DAGs in Python, cron scheduling, \cli{>>} for
  dependencies \\
\textbf{Prefect} & \cli{@flow}, \cli{@task} decorators, built-in caching
  and retries \\
\textbf{DVC pipelines} & \file{dvc.yaml}, \cli{dvc repro}, data-driven
  re-execution \\
\textbf{Feast} & \cli{feast apply}, \cli{feast materialize},
  \cli{store.get\_historical\_features()} \\
\textbf{Best practice} & Idempotent, validated, retryable, logged \\
\bottomrule
\end{tabular}
\end{cheatsheet}
